\subsection{Статические методы обнаружения аномалий}

Статические методы обнаружения аномалий являются одним из базовых подходов к обеспечению информационной безопасности и широко применяются в системах мониторинга и защиты информационных систем. Данные методы основываются на заранее заданных правилах, сигнатурах и эвристиках, определяющих допустимое поведение системы и выявляющих отклонения от него. В отличие от методов машинного обучения, статические подходы не требуют предварительного этапа обучения и функционируют на основе экспертных знаний о предметной области \cite{AD_Survey}.

В контейнерных средах, в частности в Kubernetes-кластерах, статические методы применяются для анализа audit-логов, сетевого взаимодействия, конфигурационных объектов и поведения пользователей и сервисных аккаунтов. Основной задачей является выявление заранее известных шаблонов аномального поведения, таких как попытки несанкционированного доступа, эскалации привилегий, использование запрещённых операций или отклонение от установленных политик безопасности.

К числу основных статических методов обнаружения аномалий относятся:

\begin{itemize}
    \item \textbf{Сигнатурный анализ (signature-based detection)} — основан на сопоставлении наблюдаемых событий с заранее известными шаблонами атак. Данный подход широко используется в системах обнаружения вторжений и позволяет эффективно выявлять известные угрозы с высокой точностью. Однако он не способен обнаруживать новые или модифицированные атаки \cite{Lightweight_Detection_Networks}.
    \item \textbf{Правило-ориентированные методы (rule-based detection)} — предполагают использование набора логических правил, описывающих допустимое и недопустимое поведение системы. В Kubernetes такие правила могут определять ограничения на использование привилегированных контейнеров, доступ к API или действия сервисных аккаунтов. Подобные механизмы в Kubernetes реализуются, например, с использованием Admission Controller \cite{KubernetesSecurityDocs}.
    \item \textbf{Пороговые методы (threshold-based detection)} — основаны на сравнении наблюдаемых метрик с заданными пороговыми значениями. Превышение допустимых уровней использования ресурсов (CPU, памяти) или частоты запросов может свидетельствовать о наличии аномалии. Простота реализации является преимуществом данного подхода, однако его эффективность зависит от корректности выбора порогов.
    \item \textbf{Методы белых и чёрных списков (whitelist/blacklist)} — предполагают явное задание разрешённых или запрещённых сущностей (пользователей, IP-адресов, операций). Отклонение от заданного списка интерпретируется как потенциальная аномалия. Данные методы часто применяются для контроля доступа и сетевого взаимодействия.
    \item \textbf{Методы, основанные на марковских процессах} — анализируют последовательности событий и вероятности переходов между состояниями системы.
\end{itemize}

\subsubsection{Сигнатурный анализ}

Сигнатурный анализ представляет собой метод, основанный на поиске совпадений между наблюдаемыми событиями и заранее известными шаблонами атак. Формально детектор можно представить как функцию:

$$
D(x) =
\begin{cases}
1, & \text{если } x \in S_{attack} \\
0, & \text{иначе}
\end{cases}
$$
где \(x\) — наблюдаемое событие; \\ \(S_{attack}\) — множество известных сигнатур атак \cite{IDS_survey_and_taxonomy}.

Данный подход широко применяется в системах IDS (Intrusion Detection Systems). Его основным преимуществом является высокая точность обнаружения известных атак и низкая вычислительная сложность. Однако он не способен выявлять ранее неизвестные угрозы.

\subsubsection{Правило-ориентированные методы}

Правило-ориентированные методы используют набор логических условий для описания допустимого поведения системы. Формально правило можно представить как предикат:

$$
R(x) = f(x_1, x_2, ..., x_n) \rightarrow \{0,1\},
$$
где \(x_i\) — признаки события; \\ \(f\) — логическая функция.

Инструменты для реализации такой логики (например, Falco \cite{falco}) реализуют сопоставление правил с событиями контейнеров и предоставляют как стандартный набор правил, так и механизм расширения этого набора пользовательскими правилами; аналогично, policy-ориентированные решения (Open Policy Agent \cite{opa}) применяют декларативные ограничения на этапе развертывания контейнера в кластере.

Преимущество правило-ориентированных методов состоит в их детерминированности и объяснимости. Однако та же детерминированность определяет и фундаментальное ограничение: правила не способно обнаружить ранее не встречавшиеся, тонкие или композиционные паттерны поведения (zero-day и эволюционирующие атаки), а также требуют постоянного ручного сопровождения - обновления и тестирования наборов правил по мере появления новых уязвимостей и изменчивости рабочих нагрузок.

\subsubsection{Пороговые методы}

Пороговые методы основаны на анализе статистических характеристик метрик системы. Пусть имеется наблюдаемая метрика \(x_t\), тогда аномалия определяется как:

$$
x_t \notin [\mu - k\sigma, \mu + k\sigma],
$$
где \(\mu\) — среднее значение; \\ \(\sigma\) — стандартное отклонение; \\ \(k\) — коэффициент чувствительности.

Также могут использоваться более сложные методы:
\begin{itemize}
    \item CUSUM (кумулятивные суммы);
    \item EWMA (экспоненциально взвешенное среднее);
\end{itemize}

Основное ограничение — чувствительность к изменению распределения данных и отсутствие учёта многомерных зависимостей.

\subsubsection{Методы и черных белых списков}

Метод белых списков основан на явном задании допустимых объектов:

$$
D(x) =
\begin{cases}
1, & x \in W \\
0, & x \notin W
\end{cases}
$$
где \(W\) — whitelist.

Метод черных списков, напротив, основан на запрете заранее известных нежелательных объектов:
\nopagebreak
$$
D(x) =
\begin{cases}
0, & x \in B \\
1, & x \notin B
\end{cases}
$$
где \(B\) — blacklist.

В Kubernetes такие методы применяются для контроля:

\begin{itemize}
    \item разрешённых/запрещенных IP-адресов;
    \item сервисных аккаунтов;
    \item допустимых образов контейнеров;
    \item фильтрации вредоносных сетевых запросов.
\end{itemize}

Недостатком является высокая статичность и необходимость регулярного обновления списков.

\subsubsection{Методы, основанные на марковских процессах}

Марковские методы рассматривают поведение системы как последовательность состояний \( S = \{s_1, s_2, \dots, s_n\} \), где переход между состояниями описывается условной вероятностью
$$
P(s_t \mid s_{t-1}).
$$

Таким образом, вероятность наблюдаемой последовательности состояний определяется как произведение вероятностей переходов:
$$
P(s_1, s_2, \dots, s_T) = \prod_{t=2}^{T} P(s_t \mid s_{t-1}).
$$

В задачах обнаружения аномалий модель обучается на нормальных последовательностях событий, формируя матрицу переходных вероятностей, отражающую типичное поведение системы. Если вероятность наблюдаемой последовательности оказывается ниже заданного порога, такая последовательность классифицируется как аномальная.

В контексте Kubernetes марковские модели применяются для анализа audit-логов, где состояния соответствуют типам API-запросов (например: \texttt{create}, \texttt{delete}, \texttt{patch}). Это позволяет выявлять нетипичные последовательности действий пользователей или сервисов. Основным преимуществом данного подхода является учёт временной структуры событий и зависимостей между последовательными действиями.

Однако ограничение классической марковской модели заключается в предположении о зависимости только от текущего состояния (свойство отсутствия памяти). Это снижает эффективность при анализе сложных атакующих сценариев, где важны долгосрочные зависимости и контекст.